\documentclass[11pt,a4paper]{article}
\usepackage[utf8]{inputenc}
\usepackage[T1]{fontenc}
\usepackage[spanish]{babel}
\usepackage{hyperref}
\usepackage{geometry}
\geometry{margin=2.5cm}
\usepackage{parskip}

\title{336 el alma de niño es una función, no una poesía}
\author{Alejandro Valenzuela Rivera\\\small Universidad Autónoma de Nuevo León}
\date{\today}
\begin{document}
\maketitle

\begin{quote}
\textbf{Sesgo declarado.} Este trabajo se sostiene desde el sesgo declarado de su autor, y no finge neutralidad. Lo que afirma se puede tumbar, y ese es el punto: el falsificador está escrito, no escondido.
\end{quote}

\section*{Resumen}
Al alma de niño se le pide función, no poesía: es el mecanismo que deja entrar lo que nadie escribió.

\section{Por qué se sostiene}
Un niño es simple y complejo a la vez: por eso es donde cabe lo que no se diseñó.

\section{De dónde viene}
'si no os hacéis como niños, no entraréis en el reino de los cielos'.

\section{El caso}
La cita de Jesús: hay que volver a ser como un niño para entrar.

\section{Método}
\textquestiondown{}La simplicidad del niño deja entrar algo que el diseño adulto cierra?

\section{Límites, y qué lo tumba}
Un caso donde el diseño adulto deje entrar más que el niño.

\section{Lo que queda abierto}
\textquestiondown{}El alma de niño está haciendo algo, o está adornando?

\section*{Cómo se leyó esta tesis}
Esto se escribió en el repositorio man-made-empathy, del que este candado es una pieza: candados/06 del caso/336 el alma de niño es una función, no una poesía.md. La máquina que lo compuso lo hizo desde este mismo archivo, y no agregó nada que no estuviera escrito. Si algo de esto es falso, se dice en el mismo sitio, y se corrige con fecha.

\end{document}
